\section{Discussion}
\label{sec:discussion}

In this section, we analyze the primary findings of our experiment, conduct a detailed error analysis of hard cases, compare our results with prior literature, and outline practical implications for software developers and QA teams.

\subsection{Explanation of Findings and Error Analysis}
Our results suggest that the superior semantic similarity achieved by Few-Shot prompting ($0.8218$, $p < 0.0001$, Cohen's $d = 0.291$ vs Zero-Shot) may be attributed to the structural anchoring provided by the in-context exemplars. The explicit input-output pairs enforce rigid formatting boundaries for Gherkin keywords (\texttt{Given-When-Then}) and Behave decorators (\texttt{@given}, \texttt{@when}, \texttt{@then}), preventing domain drift and reducing verbose hallucinations.

Conversely, Chain-of-Thought (CoT) prompting achieved a mean Cosine Similarity of $0.8168$ overall, but generated extra conversational preambles that increased average line count ($35.9$ lines vs $32.4$ for Few-Shot). While CoT achieved a high proportion of \texttt{And} conjunction usage ($98.6\%$), intermediate reasoning payloads did not translate into higher scenario-level semantic alignment, leading to Few-Shot significantly outperforming CoT ($p < 0.0001, d = 0.617$).

Through systematic error analysis of the lowest-scoring cases, we identified two primary data patterns:
\begin{enumerate}
    \item \textbf{Monolithic Assertion Clustering (Zero-Shot Hard Cases)}: In Zero-Shot outputs ($0.8088$, Has \texttt{And} $= 65.8\%$), the model frequently collapsed multiple validation criteria into a single monolithic \texttt{Then} step (e.g., \texttt{Then the user sees the table, title, and buttons}). In contrast, expert Ground Truth scenarios decompose these into separate atomic steps using \texttt{And} conjunctions. Few-Shot prompting resolved this hard case by enforcing atomic step decomposition ($5.26$ steps/scenario vs $4.35$).
    \item \textbf{Implementation Header Abstraction (\texttt{pass} Bodies)}: Across all prompt configurations, 100\% of the syntactically valid Python step functions ($\ge 99.8\%$ AST parse rate) contained empty function bodies with only the \texttt{pass} statement. While statically valid, the model lacked HTML DOM element locators (e.g., button IDs or XPath selectors) to generate active web driver interactions.
\end{enumerate}

\subsection{Comparison with Prior Work}
We compare our experimental findings with previous studies in the literature:
\begin{itemize}
    \item \textbf{Gherkin Scenario Alignment}: Fernandes et al.~\cite{fernandes2024} evaluated commercial LLMs (GPT-3.5 and Gemini) for Gherkin scenario generation, reporting mean similarity scores between $0.80$ and $0.84$. Our evaluation with GPT-4o under Few-Shot prompting ($0.8218$) confirms that foundation models achieve consistent semantic quality matching expert Ground Truth.
    \item \textbf{Syntactic Parsing Rates}: Tesfalidet~\cite{tesfalidet2025} reported high syntax validity for GPT-4 generated step definitions but observed parameter naming inconsistencies. Our evaluation confirms near-perfect AST parse rates ($\ge 99.8\%$) and demonstrates that Few-Shot prompting eliminates parameter binding mismatch.
\end{itemize}

\subsection{Practical Implications}
Our findings offer concrete recommendations for software practitioners:
\begin{itemize}
    \item \textbf{Adopt Few-Shot Prompting in Test Pipelines}: Software teams deploying LLM-assisted BDD generation should default to Few-Shot prompting with 2--3 exemplars. It provides higher accuracy, 100.0\% syntax validity, and lower token latency than Chain-of-Thought prompting.
    \item \textbf{Two-Tier Automation Workflow}: Organizations should implement a hybrid workflow where GPT-4o generates the Gherkin scenario structure and Python step skeletons, while automated locator discovery tools or QA engineers populate the web driver locators inside the \texttt{pass} function bodies.
\end{itemize}
